\documentclass[aspectratio=169]{beamer}

\usetheme{Madrid}
\usecolortheme{default}
\usefonttheme{serif}
\setbeamertemplate{navigation symbols}{}

\usepackage{iftex}

\ifPDFTeX
  \usepackage{newtxtext}
  \usepackage{newtxmath}
\else
  \usepackage{fontspec}
  \setmainfont{Times New Roman}
  \setsansfont{Times New Roman}
  \setmonofont{Times New Roman}
\fi

\setbeamerfont{normal text}{family=\rmfamily,series=\mdseries}
\setbeamerfont{structure}{family=\rmfamily,series=\mdseries}
\setbeamerfont{frametitle}{family=\rmfamily,series=\mdseries}
\setbeamerfont{title}{family=\rmfamily,series=\mdseries}
\setbeamerfont{itemize/enumerate body}{family=\rmfamily,series=\mdseries}
\setbeamerfont{itemize/enumerate subbody}{family=\rmfamily,series=\mdseries}
\linespread{0.96}

\begin{document}

\begin{frame}{Old Lane-Center Reference Method}
\small
\begin{columns}[T,onlytextwidth]
\begin{column}{0.48\textwidth}
\textbf{How the old method worked}
\begin{itemize}
\setlength\itemsep{0.45em}
\item The MPC reference path was built directly from the ego car to the blue dot.
\item If the ego and blue dot were in different lanes, the reference jumped to the blue-dot lane immediately.
\item There was no blending between the current lane and the target lane.
\item So the controller tried to follow the blue dot path right away, even in the middle of a junction.
\end{itemize}
\end{column}

\begin{column}{0.48\textwidth}
\textbf{Why it did not work well}
\begin{itemize}
\setlength\itemsep{0.45em}
\item At intersections there can be many nearby branches that look valid.
\item The reference path could lock onto a nearby branch instead of the branch the route really wanted.
\item A sudden jump from one lane to another made the path too aggressive.
\item As a result, the ego sometimes curved, hesitated, or pointed toward the wrong branch for a moment.
\end{itemize}
\end{column}
\end{columns}
\end{frame}

\begin{frame}{New Route-Based Lane-Center Reference}
\small
\textbf{Step-by-step idea}
\begin{itemize}
\setlength\itemsep{0.55em}
\item Step 1: keep the lane-center reference tied to the global route instead of drawing it directly from ego to blue dot.
\item Step 2: find where the ego is on the current route.
\item Step 3: build the MPC reference from the route points ahead of the ego.
\item Step 4: if a lane change is needed, move toward the new lane through the route path instead of jumping to a different branch.
\item Step 5: MPC follows this route-based reference, so the ego keeps a smoother and more stable path.
\end{itemize}

\vspace{0.6em}
\textbf{Why this works better}
\begin{itemize}
\setlength\itemsep{0.45em}
\item The reference path stays consistent through the junction.
\item The ego is less likely to bend toward a nearby but wrong branch.
\item The path remains smoother when road-boundary or lane-boundary costs are active.
\end{itemize}
\end{frame}

\begin{frame}{Extra Logic for Correct Branch Selection}
\small
\textbf{Blue-dot protection at intersections}
\begin{itemize}
\setlength\itemsep{0.55em}
\item When the blue dot looks for the next waypoint and there are multiple choices, the system first checks which branch belongs to the global route.
\item If the blue dot is inside a junction but ends up on a waypoint that is not on the route, it snaps back to the closest valid point on the route ahead.
\item After that, the blue dot continues rolling along the correct route branch.
\end{itemize}

\vspace{0.6em}
\textbf{Extra logic used for the ego path}
\begin{itemize}
\setlength\itemsep{0.55em}
\item The MPC reference is kept on the same branch that the blue dot is following.
\item The reference is built with continuity, so it does not suddenly jump to another nearby branch for one step.
\item Together, these checks make the ego follow the same branch as the blue dot through the intersection.
\end{itemize}
\end{frame}

\end{document}
